\begin{frame}{Einsatzebenen}
  \textbf{Non-Photorealistic Rendering (NPR)} kann auf drei Ebenen eingesetzt werden:
  \begin{itemize}
    \item Geometrieraum
    \item Projektionsraum
    \item Bildraum
  \end{itemize}
\end{frame}

\subsection{Techniken im Geometrieraum}
\begin{frame}{Techniken im Geometrieraum}
  \begin{itemize}
    \item direkter Einfluss auf die 3D-Geometrie der Objekte
    \item Nachteil: erfordert einen hohen Zeitaufwand
    \item wird daher selten bevorzugt
    \item Techniken:
    \begin{itemize}
      \item Computer Sketching
      \item Noise Modifier
      \item Level of Detail
    \end{itemize}
  \end{itemize}
\end{frame}

\begin{frame}{Computer Sketching}
  \begin{itemize}
    \item Paul Bourk, 1990
    \item Frage: Können CAD-Anwendungen auch unscharfe,
    an menschliche Skizzen erinnernde Entwürfe
    ausgeben?
    \item Kriterien für Skizzen:
    \begin{itemize}
      \item Linien enden nicht exakt an Schnittpunkten (Kreuzungen)
      \item Linien sind nicht gradlinig sondern "`verwackelt"'
    \end{itemize}
  \end{itemize}
\end{frame}

\begin{frame}{Computer Sketching}
  \begin{itemize}
    \item Hauptansatz: direkte Manipulation der 3D-Geometrie
    \item Umsetzung:
    \begin{itemize}
      \item Linien werden um randomisierte Faktoren verlängert
      \item Aufspalten der Linien und Mittelpunktverschiebung
    \end{itemize}
  \end{itemize}
\end{frame}

\begin{frame}{Computer Sketching}
  \includegraphics[width=0.5\textwidth]{../images/haus.jpg}
\end{frame}

\begin{frame}{Noise Modifier}
  \begin{itemize}
    \item "`Rauschen"' in 3D-Modellen
    \item Geometriepunkte werden um einen zufälligen Faktor
    verschoben
    \item meist Anwendung auf glatte Flächen, um
    "`natürliche"' Flächen zu erzeugen (Landschaften)
  \end{itemize}
\end{frame}

\begin{frame}{Level of Detail}
  \begin{itemize}
    \item Einsatz in vielen Bereichen:
    \begin{itemize}
      \item Kartographie: angezeigte Detailstufe in Abhängigkeit vom
      Maßstab (Generalisierung)
      \item photorealistische Darstellungen: angezeigte Detailstufe der
      Szene in Abhängigkeit von der Betrachterposition
      \item NPR: Detailstufe in Abhängigkeit zur "`Wichtigkeit"' der
      Objekte
      \begin{itemize}
        \item Eine Möglichkeit ist die Meshsimplifizierung
      \end{itemize}
    \end{itemize}
  \end{itemize}
\end{frame}

\begin{frame}{Level of Detail}
  \begin{itemize}
    \item Meshsimplifizierung:
    \begin{itemize}
      \item reduziert die Anzahl der Polygone eines 3D Modells
      \item äußere Form wird dabei approximiert
      \item Dadurch: schnellere Verarbeitung in Echtzeitsystemen möglich
    \end{itemize}
  \end{itemize}
\end{frame}

\subsection{Techniken im Projektionsraum}
\begin{frame}{Techniken im Projektionsraum}
  \begin{itemize}
    \item meistens Kombination aus Projektions- und
    Bildraumverfahren
    \item dadurch leichteres Arbeiten mit 3D-Modellen und
    Erhöhung der grafischen Qualität der Ausgabe
    \item Techniken:
    \begin{itemize}
      \item G-Buffer
      \item Line-Rendering
      \item Stroke Textures
    \end{itemize}
  \end{itemize}
\end{frame}

\begin{frame}{G-Buffer}
  \begin{itemize}
    \item "`Geometrie-Buffer"' Saito \& Takahashi, 1990
    \item in den meisten kommerziellen Cartoon-Render-Paketen implementiert
    \item Prinzip: zusätzliche Speicherung von geometrischen
    Daten zur Informationsgewinnung im Bildraum
    \begin{itemize}
      \item Flächen-Normalen des 3D-Modells
      \item UV-Koordinaten des 3D-Modells
      \item Z-Buffer des 3D-Modells
    \end{itemize}
    \item Grundidee: Informationen zur Bilderzeugung aus dem G-Buffer
    extrahieren, analysieren und zusammenfügen.  
  \end{itemize}
\end{frame}

\begin{frame}{Line-Rendering}
  \begin{itemize}
    \item Simulation der Semantik von Linien
    \item geometrische Semantik:
    \begin{itemize}
      \item Begrenzungen, Silhouetten, Diskontinuitäten (Falten)
    \end{itemize}
    \item stilistische Semantik:
    \begin{itemize}
      \item Stärke der Linien, Transparenz, Linientyp (durchgezogen,
      gestrichelt, etc.)
    \end{itemize}
  \end{itemize}
\end{frame}

\begin{frame}{Line-Rendering}
  \begin{itemize}
    \item Abspeicherung dieser zusätzlichen Attribute in
    geeigneter Form
    \begin{itemize}
      \item Eine Möglichkeit: Speicherung in Matrixform
    \end{itemize}
    \item Generierung der Attribute
    \begin{itemize}
      \item z.B. durch Inferenzregeln zu Bestimmung der
      Charakteristika der einzelnen Linien
    \end{itemize}
    \item Gute Kombinationsmöglichkeiten mit anderen
    NPR-Techniken (etwa G-Buffer)
  \end{itemize}
\end{frame}

\begin{frame}{Stroke Textures}
  \begin{itemize}
    \item Winkenbach \& Salesin, 1994
    \item Prinzip basiert auf konventionellen 2D-Malsystem
    \begin{itemize}
      \item "`strichbasierte"' Texturen werden angefertigt
      \item mittels Texturemapping auf 3D-Objekte aufgebracht
      \item Dabei: Anwendung des LOD-Prinzips auf die Textur
    \end{itemize}
  \end{itemize}
\end{frame}

\begin{frame}{Stroke Textures}
  \includegraphics[width=\textwidth]{../images/strokes.jpg}
  \begin{itemize}
    \item Linienstärken und "`Detail Tags"' sind durch den
    Benutzer spezifizierbar
  \end{itemize}
\end{frame}

\subsection{Techniken im Bildraum}
\begin{frame}{Techniken im Bildraum}
  \begin{itemize}
    \item vorhandenes Bild wird analysiert und neu
    gezeichnet
    \item Techniken:
    \begin{itemize}
      \item Texture Elements
      \item Hairy brushes
      \item Iterated Function System (IFS)
    \end{itemize}
  \end{itemize}
\end{frame}

\begin{frame}{Texture Elements}
  \begin{itemize}
    \item Mezei et al., 1974
    \item Ursprünglich: Synthese sehr "`naturalistischer"' Bilder
    \item Beobachtung: in der Natur kommen sich oft
    wiederholende Muster vor
  \end{itemize}
\end{frame}

\begin{frame}{Texture Elements}
  \begin{itemize}
    \item Umsetzung: vorgefertigte grafische Subelemente
    werden zusammengefügt, um neue Bilder zu
    erzeugen
    \item Variationen durch Skalierung, Rotation und andere
    Verzerrungen
    \item Erzielt werden Verallgemeinerungen der
    Oberflächeneigenschaften, die wie menschliche
    Illustrationen wirken
  \end{itemize}
\end{frame}

\begin{frame}{Hairy Brushes}
  \begin{itemize}
    \item Strassmann, 1986
    \item Verhalten von Pinsel und Farbe soll simuliert werden
    \item Ansatz: Pinsel als 1D-Abdruck der Pinselborsten
    \item Borsten laufen an einer parametrischen Kurve
    entlang, die durch Knoten definiert ist
  \end{itemize}
\end{frame}

\begin{frame}{Hairy Brushes}
  \begin{itemize}
    \item Knoten enthalten Informationen über Position und
    Druck des Pinsels
    \item jeder "`Strich"' wird mit einer speziellen Menge an
    Farbe gezeichnet
    \item Menge nimmt mit der Länge der Linie ab
    \item erneutes "`Eintunken"' des Pinsels ist möglich
  \end{itemize}
\end{frame}

\begin{frame}{Iterated Function Systems (IFS)}
\begin{itemize}
 \item IFS: Systeme von Funktionen die wiederholt
 ausgeführt werden
 \item Einsatz: zur Beschreibung von Fraktalen
 \item Fraktale sind selbstbezüglich (rekursiv)
 \begin{itemize}
   \item Beispiel: Barnsley-Farn
     \begin{columns}
    \column{.2\textwidth}
	\includegraphics[width=\textwidth]{../images/BarnsleyFarn.jpg}
	\column{.2\textwidth}
    \includegraphics[width=\textwidth]{../images/BarnsleyFarnDetail.jpg}
  \end{columns}
 \end{itemize}
\end{itemize}

\end{frame}

\begin{frame}{Iterated Function Systems (IFS)}
  \begin{itemize}
    \item Nutzung des Prinzips zur fraktalen Kompression von
    Bildern (Barnsley, 1988)
    \item Kompressionstechnik basiert auf dem "`Collage-
    Theorem"' (Selbstähnlichkeit von Bildern)
  \end{itemize}
\end{frame}